\chapter{Prawa wielkich liczb}
\label{ch:lln}

\providecommand{\atlasref}[2]{#2}

\paragraph{Odsyłacze międzytomowe.}
Zastosowania mocnego prawa wielkich liczb do zgodności
estymatorów i metod Monte Carlo przedstawiono w tomie IV
(w rozdziałach o estymatorach oraz metodach Monte Carlo%
\IfFileExists{probability_volume_iv.aux}{, odpowiednio~\ref{IV-ch:estimators} i~\ref{IV-ch:monte-carlo}}{}).


Prawa wielkich liczb formalizują stabilizację częstości
obserwowanych w długiej serii doświadczeń. W serii rzutów symetryczną monetą
względna liczba orłów nie musi być bliska $1/2$ po
kilkunastu próbach, lecz przy dużej liczbie niezależnych
rzutów odchylenia stają się rzadkie. Zależnie od tego,
czy badamy prawdopodobieństwo odchylenia dla dużych $n$,
czy całe nieskończone trajektorie średnich, otrzymujemy
słabe albo mocne prawo wielkich liczb.

\begin{lemat}[Tożsamość warstwowa dla wartości oczekiwanej]
Niech $Y$ będzie nieujemną rzeczywistą zmienną losową.
Wówczas, dopuszczając wartość nieskończoną,
\[
\E Y=\int_0^\infty \Prob(Y>t)\,\dd\lambda_1(t).
\]
W szczególności, gdy $\E Y<\infty$, to
$\sum_{k=1}^\infty\Prob(Y>k)\leq\E Y$.
\end{lemat}
\begin{proof}
Pierwsza równość jest tożsamością ogonową
wykazaną w rozdziale~\ref{ch:distributions}
(w przykładzie porównującym ogony rozkładów
normalnego i Cauchy'ego); jej dowód z twierdzenia
Tonellego dotyczy dowolnej nieujemnej zmiennej
losowej, niezależnie od skończoności jej średniej.
Funkcja ogonowa jest mierzalna, gdyż jest nierosnąca.
Dla każdego całkowitego $k\geq1$ i $t\in[k-1,k)$
z inkluzji $\{Y>k\}\subseteq\{Y>t\}$ wynika
$\Prob(Y>t)\geq\Prob(Y>k)$, zatem
\[
\sum_{k=1}^N\Prob(Y>k)
\leq\sum_{k=1}^N\int_{k-1}^k\Prob(Y>t)\,\dd\lambda_1(t)
\leq\E Y.
\]
Przejście z $N\to\infty$ dowodzi nierówności
dla szeregu.
\end{proof}

\section{Sformułowania i wersja wariancyjna}

\begin{definicja}
Niech $(X_n)_{n\geq1}$ będzie ciągiem rzeczywistych
zmiennych losowych na jednej przestrzeni
probabilistycznej i niech
$\overline X_n=n^{-1}\sum_{k=1}^nX_k$ dla $n\geq1$.
Każda średnia jest mierzalna jako skończona kombinacja
liniowa zmiennych mierzalnych. Jeżeli wszystkie $X_k$
są całkowalne, to także każda średnia jest całkowalna.
Sama definicja zbieżności średnich nie wymaga jednak
istnienia ich wartości oczekiwanych.
Dla $m\in\R$ mówimy, że zachodzi \emph{słabe prawo
wielkich liczb} (SPWL) z granicą $m$, jeśli
$\overline X_n\xrightarrow{\Prob}m$.
Mówimy, że zachodzi \emph{mocne prawo wielkich liczb}
(MPWL) z granicą $m$, jeśli
$\overline X_n\to m$ prawie na pewno.
\end{definicja}

Mocne prawo pociąga słabe, ponieważ zbieżność prawie na
pewno implikuje zbieżność według prawdopodobieństwa
(rozdział~\ref{ch:convergence}).
W drugą stronę implikacja nie zachodzi dla dowolnych ciągów
zmiennych losowych. Istotnie, niech $(A_n)_{n\geq1}$
będzie ciągiem funkcji charakterystycznych przedziałów
diadycznych z przykładu w rozdziale~\ref{ch:convergence},
zbieżnym do zera według prawdopodobieństwa, lecz
niezbieżnym punktowo w żadnym punkcie $[0,1)$.
Połóżmy $A_0=0$ i $X_n=nA_n-(n-1)A_{n-1}$.
Każda zmienna $X_n$ jest mierzalna i ograniczona,
a sumowanie teleskopowe daje
\[
\overline X_n=\frac1n\sum_{k=1}^n
(kA_k-(k-1)A_{k-1})=A_n.
\]
Dla tego ciągu zachodzi SPWL z granicą zero, ale nie MPWL.
Nawet przy zbieżności prawie na pewno nie twierdzimy,
że średnia od pewnego momentu jest równa $m$.
Twierdzimy, że poza pewnym mierzalnym zbiorem zerowym
dla każdego $\varepsilon>0$ istnieje indeks
$N=N(\omega,\varepsilon)$, po którym
$|\overline X_n(\omega)-m|<\varepsilon$ dla wszystkich
$n\geq N$.

\begin{twierdzenie}[SPWL z kontroli wariancji]
Niech $(X_n)_{n\geq1}$ będzie ciągiem rzeczywistych zmiennych
losowych na wspólnej przestrzeni probabilistycznej,
o skończonych drugich momentach. Jeśli dla pewnego $m\in\R$
\[
\frac1n\sum_{k=1}^n\E X_k\longrightarrow m,
\qquad
\frac1{n^2}\Var\!\left(\sum_{k=1}^nX_k\right)\longrightarrow0,
\]
to $\overline X_n\to m$ w $L^2$, a w szczególności
$\overline X_n\xrightarrow{\Prob}m$.
W szczególności wystarcza, by $X_k$ były parami
nieskorelowane, miały wspólną średnią $m$ i
$\sum_{k=1}^n\Var(X_k)=o(n^2)$.
\end{twierdzenie}
\begin{proof}
Z nierówności Cauchy'ego--Schwarza dla $|X_k|$ i
funkcji stałej $1$ wynika
$\E|X_k|\leq(\E X_k^2)^{1/2}<\infty$.
Skończone sumy także należą do $L^2$, ponieważ
punktowo $(\sum_{k=1}^n X_k)^2\leq n\sum_{k=1}^n X_k^2$
z nierówności Cauchy'ego--Schwarza w $\R^n$.
Po scałkowaniu prawa strona jest skończona.
Oznaczmy $a_n=n^{-1}\sum_{k=1}^n\E X_k$.
Z liniowości wartości oczekiwanej
$\E\overline X_n=a_n$, zaś z definicji wariancji
\[
\Var(\overline X_n)
=\E\left|\frac1n\sum_{k=1}^n(X_k-\E X_k)\right|^2
=\frac1{n^2}\Var\left(\sum_{k=1}^nX_k\right).
\]
W szczególności z obu założeń twierdzenia wynika
\[
\E|\overline X_n-m|^2
=\Var(\overline X_n)+|a_n-m|^2\longrightarrow0,
\]
gdyż po rozwinięciu kwadratu wyraz mieszany znika
na mocy $\E(\overline X_n-a_n)=0$.
Dla ustalonego $\varepsilon>0$ i dostatecznie
dużego $n$ mamy $|a_n-m|<\varepsilon/2$.
Z nierówności trójkąta
$|\overline X_n-m|\leq|\overline X_n-a_n|+|a_n-m|$
wynika inkluzja zdarzeń
$\{|\overline X_n-m|>\varepsilon\}\subseteq
\{|\overline X_n-a_n|>\varepsilon/2\}$.
Z nierówności Czebyszewa (rozdział~\ref{ch:expectation}) otrzymujemy
\[
\Prob(|\overline X_n-m|>\varepsilon)
\leq\frac{4}{\varepsilon^2}\Var(\overline X_n)
\longrightarrow0.
\]
Rozwijając kwadrat sumy zmiennych scentrowanych
i korzystając z liniowości wartości oczekiwanej, mamy
\[
\Var\left(\sum_{k=1}^n X_k\right)
=\sum_{k=1}^n\Var(X_k)
+2\sum_{1\leq j<k\leq n}\operatorname{Cov}(X_j,X_k).
\]
Iloczyny zmiennych scentrowanych są całkowalne
z nierówności Cauchy'ego--Schwarza.
Dla zmiennych parami nieskorelowanych wszystkie
kowariancje w drugiej sumie są zerowe. Przy wspólnej
średniej $a_n=m$ oraz $\sum_{k=1}^n\Var(X_k)=o(n^2)$
obie hipotezy twierdzenia są zatem spełnione.
\end{proof}

\begin{wniosek}[SPWL dla parami niezależnych zmiennych o wspólnej skończonej wariancji]
Jeśli rzeczywiste zmienne $X_1,X_2,\ldots$ na jednej
przestrzeni probabilistycznej są parami niezależne,
mają wspólną średnią $m$ i wspólną skończoną wariancję $\sigma^2$,
to $\overline X_n\xrightarrow{\Prob}m$ oraz
$\E|\overline X_n-m|^2=\sigma^2/n\to0$.
\end{wniosek}
\begin{proof}
Z niezależności parami wynika
$\E[(X_j-m)(X_k-m)]=\E(X_j-m)\E(X_k-m)=0$ dla $j\ne k$
(rozdział~\ref{ch:expectation}). Zatem ze wzoru
na wariancję sumy w poprzednim dowodzie mamy
$\Var(\sum_{k=1}^nX_k)=n\sigma^2$.
Po podzieleniu przez $n^2$ dostajemy
$\Var(\overline X_n)=\sigma^2/n$.
Ponieważ $\E\overline X_n=m$, ostatnia wielkość
jest równa $\E|\overline X_n-m|^2$.
Oznacza to zbieżność średnich do stałej $m$ w $L^2$,
a zatem według prawdopodobieństwa
(rozdział~\ref{ch:convergence}).
\end{proof}

\section{SPWL przy samym istnieniu wartości oczekiwanej}

Oszacowanie za pomocą wariancji w dowodzie SPWL
wymaga skończonego drugiego momentu, aby dało
niebanalne ograniczenie. Istnieją jednak rozkłady
o skończonej średniej i nieskończonej wariancji,
np. rozważany dalej rozkład Pareto z parametrem
$1<\alpha<2$. Dla parami niezależnych prób o tym
samym rozkładzie wystarcza sama całkowalność.

\begin{twierdzenie}[Słabe prawo wielkich liczb w wersji całkowalnej]
Niech $X_1,X_2,\ldots$ będą rzeczywistymi zmiennymi
losowymi na jednej przestrzeni probabilistycznej,
parami niezależnymi, o tym samym rozkładzie, oraz
niech $\E|X_1|<\infty$. Wówczas, oznaczając $m=\E X_1$,
mamy $\overline X_n\xrightarrow{\Prob}m$.
\end{twierdzenie}
\begin{proof}
Dla $K>0$ obetnijmy duże wartości:
$Y_k^{(K)}=X_k\chi_{\{|X_k|\leq K\}}$.
Funkcja $g_K(x)=x\chi_{[-K,K]}(x)$ jest borelowska,
więc $Y_k^{(K)}=g_K(X_k)$ są mierzalnymi zmiennymi
losowymi. Dla $j\ne k$ i borelowskich $A,B\subseteq\R$
zdarzenia $\{Y_j^{(K)}\in A\}$ i $\{Y_k^{(K)}\in B\}$
są odpowiednio zdarzeniami
$\{X_j\in g_K^{-1}(A)\}$ i $\{X_k\in g_K^{-1}(B)\}$,
a więc są niezależne. Zmienne obcięte mają wspólny
rozkład jako obrazy tego samego rozkładu przez $g_K$.
Ponieważ $|Y_k^{(K)}|\leq K$, ich drugie momenty są
nie większe niż $K^2$. Poprzedni wniosek ma więc zastosowanie z
$m_K=\E Y_1^{(K)}$ i $\sigma_K^2=\Var(Y_1^{(K)})$.
Dla ustalonego $K$ otrzymujemy
\[
\E\left|\frac1n\sum_{k=1}^nY_k^{(K)}-m_K\right|^2
=\frac{\sigma_K^2}{n}\leq\frac{K^2}{n}.
\]
W szczególności $n^{-1}\sum_{k=1}^nY_k^{(K)}
\xrightarrow{\Prob}m_K$.

Kontrolujemy teraz błąd obcięcia, bez zakładania
skończonej wariancji $X_1$. Każda zmienna $X_k$ jest
całkowalna ze wzoru \eqref{eq:lotus} i jednakowości
rozkładów. Ponieważ
$|X_k-Y_k^{(K)}|=|X_k|\chi_{\{|X_k|>K\}}$,
liniowość wartości oczekiwanej, nierówność trójkąta
oraz jednakowość rozkładów dają
\[
\E\left|\frac1n\sum_{k=1}^n(X_k-Y_k^{(K)})\right|
\leq\frac1n\sum_{k=1}^n\E|X_k-Y_k^{(K)}|
=\E[|X_1|\chi_{\{|X_1|>K\}}]=:u(K).
\]
Dla całkowitych $K\to\infty$ funkcje
$|X_1|\chi_{\{|X_1|>K\}}$ zbiegają punktowo do zera
i są zmajoryzowane przez całkowalną funkcję $|X_1|$.
Z \atlasref{ME-08::thm:dct}{twierdzenia o zbieżności zmajoryzowanej}
wynika $u(K)\to0$; dla dowolnych rzeczywistych $K\to\infty$
wniosek wynika z monotoniczności $u$.
Ponadto
$|\E Y_1^{(K)}-m|\leq\E|Y_1^{(K)}-X_1|=u(K)$.
Ustalmy $\varepsilon>0$ i wybierzmy $K$ takie,
że $u(K)<\varepsilon/3$.
Oznaczmy $\overline Y_n^{(K)}=n^{-1}\sum_{k=1}^nY_k^{(K)}$.
Mamy dokładną tożsamość
\[
\overline X_n-m
=(\overline X_n-\overline Y_n^{(K)})
+(\overline Y_n^{(K)}-\E Y_1^{(K)})
+(\E Y_1^{(K)}-m).
\]
Dla modułu pierwszego składnika stosujemy
nierówność Markowa (rozdział~\ref{ch:expectation})
oraz wcześniej uzyskane oszacowanie wartości oczekiwanej:
\[
\Prob\!\left(\left|\frac1n\sum_{k\leq n}
(X_k-Y_k^{(K)})\right|>\varepsilon/3\right)
\leq\frac3\varepsilon\E|X_1-Y_1^{(K)}|.
\]
Trzeci składnik ma moduł mniejszy niż $\varepsilon/3$.
Z nierówności trójkąta i podaddytywności prawdopodobieństwa
otrzymujemy
\[
\begin{aligned}
\Prob(|\overline X_n-m|>\varepsilon)
&\leq\frac{3u(K)}{\varepsilon}
+\Prob\left(\left|\frac1n\sum_{k=1}^nY_k^{(K)}
-\E Y_1^{(K)}\right|>\frac{\varepsilon}{3}\right).
\end{aligned}
\]
Dla ustalonego $K$ drugi składnik dąży do zera
na mocy wersji wariancyjnej SPWL.
Zatem dla wszystkich dostatecznie dużych $K$
\[
0\leq\limsup_{n\to\infty}
\Prob(|\overline X_n-m|>\varepsilon)
\leq\frac{3u(K)}{\varepsilon}.
\]
Przechodząc z $K\to\infty$, otrzymujemy tezę.
\end{proof}

\section{Nierówność maksymalna Kołmogorowa}

W dowodzie mocnego prawa sama zbieżność do zera
prawdopodobieństw odchyleń dla pojedynczych indeksów
nie wystarcza. Potrzebna jest dodatkowa kontrola
zdarzeń odchyleń na całym ogonie ciągu.
Nierówność maksymalna pozwala przejść od sumy wariancji
do prawdopodobieństwa, że \emph{któraś} suma częściowa
przekroczy próg.

\begin{twierdzenie}[Nierówność maksymalna Kołmogorowa]
Niech $n\geq1$ będzie liczbą całkowitą, a
$Z_1,\ldots,Z_n$ będą wzajemnie niezależnymi rzeczywistymi
zmiennymi losowymi na jednej przestrzeni probabilistycznej
$(\Omega,\mathcal F,\Prob)$,
o $\E Z_k=0$ i $\E Z_k^2<\infty$.
Dla $S_j=\sum_{k=1}^j Z_k$ oraz $a>0$ zachodzi
\[
\Prob\!\left(\max_{1\leq j\leq n}|S_j|\geq a\right)
\leq\frac1{a^2}\sum_{k=1}^n\Var(Z_k).
\]
\end{twierdzenie}
\begin{proof}
Wszystkie sumy częściowe należą do $L^2$, a ich
średnie wynoszą zero z liniowości wartości oczekiwanej.
Dla $1\leq k\leq n$ niech
\[
A_k=\left(\bigcap_{j=1}^{k-1}\{|S_j|<a\}\right)
\cap\{|S_k|\geq a\},
\]
przy czym dla $k=1$ puste przecięcie oznacza całą
przestrzeń $\Omega$. Są to mierzalne zdarzenia pierwszego
osiągnięcia progu. Jeżeli $k<\ell$, to na $A_k$
mamy $|S_k|\geq a$, podczas gdy na $A_\ell$ mamy
$|S_k|<a$, stąd zdarzenia są parami rozłączne.
Każdy wynik, dla którego maksimum osiąga próg,
należy do $A_k$ o najmniejszym indeksie osiągnięcia,
a każdy wynik z $A_k$ spełnia warunek na maksimum.
Zatem
\[
\bigcup_{k=1}^n A_k
=\left\{\max_{1\leq j\leq n}|S_j|\geq a\right\}.
\]

Dla $k<n$ oznaczmy $U=(Z_1,\ldots,Z_k)$ oraz
$V=(Z_{k+1},\ldots,Z_n)$.
Jeżeli $\mu_i=\Law(Z_i)$, to z charakteryzacji wzajemnej
niezależności przez iloczynowy rozkład łączny
(rozdział~\ref{ch:random-variables}) oraz
\atlasref{ME-12::cor:associative-product-12}{łączności konstrukcji miary produktowej}
wynika
\[
\Law(U,V)
=\bigotimes_{i=1}^n\mu_i
=\left(\bigotimes_{i=1}^k\mu_i\right)
 \otimes\left(\bigotimes_{i=k+1}^n\mu_i\right)
=\Law(U)\otimes\Law(V).
\]
Przyjmujemy tu naturalne utożsamienie
$\R^k\times\R^{n-k}$ z $\R^n$.
Zmienna $S_k\chi_{A_k}$ jest borelowską funkcją
pierwszego wektora, ponieważ warunki definiujące
$A_k$ są przeciwobrazami przedziałów przez ciągłe
funkcje sum częściowych. Zmienna $S_n-S_k$ jest
ciągłą funkcją drugiego wektora. Dla zbiorów
borelowskich $B,C\subseteq\R$ zdarzenie
$\{S_k\chi_{A_k}\in B,\ S_n-S_k\in C\}$ jest
przeciwobrazem odpowiedniego prostokąta przez $(U,V)$.
Jeśli $f(U)=S_k\chi_{A_k}$ i $g(V)=S_n-S_k$, to
\[
\begin{aligned}
\Prob(f(U)\in B,\ g(V)\in C)
&=(\Law(U)\otimes\Law(V))(f^{-1}(B)\times g^{-1}(C))\\
&=\Prob(f(U)\in B)\Prob(g(V)\in C).
\end{aligned}
\]
Dowodzi to niezależności obu zmiennych.
Dla $k=n$ druga zmienna jest identycznie zerowa,
więc niezależność zachodzi także w tym przypadku.
Ponadto $\E(S_n-S_k)=0$.
Wszystkie wyrazy w poniższym rozwinięciu są całkowalne:
$S_k\chi_{A_k}$ oraz $S_n-S_k$ należą do $L^2$,
a ich iloczyn należy do $L^1$ z nierówności
Cauchy'ego--Schwarza. Rozwijamy kwadrat:
\[
\E[S_n^2\chi_{A_k}]
=\E[S_k^2\chi_{A_k}]
+2\E[S_k(S_n-S_k)\chi_{A_k}]
+\E[(S_n-S_k)^2\chi_{A_k}].
\]
Z faktoryzacji wartości oczekiwanej iloczynu niezależnych
zmiennych (rozdział~\ref{ch:expectation}) otrzymujemy
\[
\E[S_k\chi_{A_k}(S_n-S_k)]
=\E[S_k\chi_{A_k}]\E(S_n-S_k)=0.
\]
Na $A_k$ zachodzi $S_k^2\geq a^2$, a poza $A_k$
obie strony nierówności
$S_k^2\chi_{A_k}\geq a^2\chi_{A_k}$ są zerowe.
Pierwsza całka jest więc nie mniejsza niż
$a^2\Prob(A_k)$, a ostatnia jest nieujemna.
Sumując po $k$ i wykorzystując rozłączność zdarzeń,
otrzymujemy
\[
\begin{aligned}
\E S_n^2
&\geq \E\left[S_n^2\chi_{\bigcup_{k=1}^n A_k}\right]
=\sum_{k=1}^n\E[S_n^2\chi_{A_k}]\\
&\geq a^2\sum_{k=1}^n\Prob(A_k)
=a^2\Prob\left(\max_{1\leq j\leq n}|S_j|\geq a\right).
\end{aligned}
\]
Ponieważ $\E S_n=0$ i kowariancje różnych składników
są zerowe, z wcześniej wyprowadzonego wzoru na
wariancję sumy mamy
$\E S_n^2=\Var(S_n)=\sum_{k=1}^n\Var(Z_k)$.
Dzielenie przez $a^2>0$ kończy dowód.
\end{proof}

\begin{lemat}[Kryterium zbieżności szeregu niezależnych zmiennych]
Niech $(Z_n)_{n\geq1}$ będzie ciągiem wzajemnie
niezależnych rzeczywistych zmiennych losowych na jednej
przestrzeni probabilistycznej, o zerowych średnich
i skończonych wariancjach. Jeśli $\sum_{n=1}^\infty\Var(Z_n)<\infty$,
to szereg $\sum_{n=1}^\infty Z_n$ jest zbieżny
prawie na pewno.
\end{lemat}
\begin{proof}
Niech $S_m=\sum_{k=1}^m Z_k$, $S_0=0$.
Z założenia ogony szeregu nieujemnych wariancji
dążą do zera. Możemy więc rekurencyjnie wybrać
ściśle rosnące indeksy $n_j$ dla $j\geq1$ tak, by
$\sum_{k>n_j}\Var(Z_k)\leq2^{-3j}$.
Zastosujmy nierówność maksymalną do skończonego
ciągu $Z_{n_j+1},\ldots,Z_m$. Dla każdego $m>n_j$
\[
\Prob\!\left(\max_{n_j<r\leq m}
\left|\sum_{k=n_j+1}^rZ_k\right|>2^{-j}\right)
\leq 2^{2j}\sum_{k>n_j}\Var(Z_k)\leq2^{-j}.
\]
Dla ustalonego $j$ zdarzenia po lewej stronie
rosną wraz z $m$ do mierzalnego zdarzenia
\[
B_j=\bigcup_{r>n_j}\{|S_r-S_{n_j}|>2^{-j}\}
=\left\{\sup_{r>n_j}|S_r-S_{n_j}|>2^{-j}\right\}.
\]
W ostatniej równości ważna jest ostra nierówność.
Z ciągłości miary z dołu (rozdział~\ref{ch:probability})
mamy $\Prob(B_j)\leq2^{-j}$.
Ponieważ $\sum_{j\geq1}\Prob(B_j)\leq1$, pierwszy
lemat Borela--Cantellego (rozdział~\ref{ch:probability})
daje $\Prob(\limsup_j B_j)=0$.
Dla $\omega$ spoza tego mierzalnego zbioru zerowego
istnieje $J(\omega)$ takie, że dla każdego $j\geq J(\omega)$
i każdego $r>n_j$ zachodzi
$|S_r(\omega)-S_{n_j}(\omega)|\leq2^{-j}$.
Dla dowolnych $r,s\geq n_j$ nierówność trójkąta daje
$|S_r(\omega)-S_s(\omega)|\leq2^{1-j}$.
W przypadku indeksu równego $n_j$ odpowiadające
odchylenie od $S_{n_j}$ jest zerowe, więc oszacowanie
obejmuje również ten przypadek.
Dla dowolnego $\eta>0$ wybieramy $j\geq J(\omega)$
takie, że $2^{1-j}<\eta$, i przyjmujemy $n_j$ jako
indeks w warunku Cauchy'ego. Sumy częściowe są więc
ciągiem Cauchy'ego dla każdej takiej $\omega$. Zupełność $\R$
zapewnia ich zbieżność do skończonej granicy.
Oznaczmy $G=\Omega\setminus\limsup_j B_j$.
Funkcje $S_m\chi_G$ są mierzalne i zbiegają punktowo
na całej przestrzeni: na $G$ do skończonej granicy
sum częściowych, a na $G^c$ do zera. Ich granica
jest więc mierzalną rzeczywistą zmienną losową.
Szereg zbiega więc prawie na pewno.
\end{proof}

\begin{lemat}[Lemat Kroneckera]
Niech $(a_n)_{n\geq1}$ będzie ciągiem liczb rzeczywistych.
Jeżeli szereg $\sum_{n=1}^\infty a_n/n$ jest zbieżny,
to $n^{-1}\sum_{k=1}^n a_k\to0$.
\end{lemat}
\begin{proof}
Niech $s_n=\sum_{k=1}^n a_k/k$ oraz $s_0=0$.
Z założenia $s_n\to s\in\R$.
Ponieważ $a_k=k(s_k-s_{k-1})$, otrzymujemy
\[
\begin{aligned}
\sum_{k=1}^n a_k
&=\sum_{k=1}^n k s_k-\sum_{k=0}^{n-1}(k+1)s_k\\
&=n s_n-\sum_{k=1}^{n-1}s_k,
\end{aligned}
\]
gdzie użyliśmy $s_0=0$.
Dzieląc przez $n$, otrzymujemy
$\frac1n\sum_{k=1}^n a_k
=s_n-\frac1n\sum_{k=1}^{n-1}s_k$.
Pierwszy składnik dąży do $s$. Dla drugiego zauważmy,
że
\[
\left|\frac1n\sum_{k=1}^{n-1}s_k-s\right|
\leq\frac1n\sum_{k=1}^{n-1}|s_k-s|+\frac{|s|}{n}.
\]
Dla dowolnego $\eta>0$ wybieramy $K\geq1$ takie,
że $|s_k-s|<\eta$ dla $k\geq K$.
Dla $n>K$ prawa strona jest nie większa niż
$n^{-1}\sum_{k=1}^{K-1}|s_k-s|+\eta+|s|/n$.
Jej granica górna po $n$ nie przekracza $\eta$,
a $\eta$ jest dowolne. Zatem drugi składnik także
dąży do $s$, co dowodzi tezy. Jest to zastosowanie
twierdzenia Cesàra o średnich ciągu zbieżnego.
\end{proof}

\begin{twierdzenie}[MPWL przy sumowalnych wariancjach ważonych]
Niech $X_1,X_2,\ldots$ będą wzajemnie niezależnymi
rzeczywistymi zmiennymi losowymi na jednej przestrzeni
probabilistycznej, o skończonych drugich momentach.
Jeśli
$\sum_{n=1}^\infty\Var(X_n)/n^2<\infty$,
to
\[
\frac1n\sum_{k=1}^n(X_k-\E X_k)\longrightarrow0
\quad\text{prawie na pewno}.
\]
Jeżeli ponadto dla pewnego $m\in\R$ zachodzi
$n^{-1}\sum_{k=1}^n\E X_k\to m$,
to $\overline X_n\to m$ prawie na pewno.
\end{twierdzenie}
\begin{proof}
Zmienna $Z_n=(X_n-\E X_n)/n$ ma średnią zero
i wariancję $\Var(X_n)/n^2$ z liniowości wartości
oczekiwanej i jednorodności wariancji.
Zmienne $Z_n$ są wzajemnie niezależne, ponieważ
każda z nich jest borelowską funkcją wyłącznie
odpowiadającego jej $X_n$ (w definicji niezależności
zastępujemy zbiory borelowskie ich przeciwobrazami).
Z założenia szereg wariancji jest zbieżny, więc na
mocy kryterium zbieżności szeregu niezależnych zmiennych
istnieje mierzalny zbiór $G$ o prawdopodobieństwie
jeden, na którym szereg $\sum_{n=1}^\infty Z_n(\omega)$
jest zbieżny. Dla każdego ustalonego $\omega\in G$
stosujemy lemat Kroneckera do liczb
$a_n=X_n(\omega)-\E X_n$. Otrzymujemy
\[
\frac1n\sum_{k=1}^n(X_k(\omega)-\E X_k)\longrightarrow0.
\]
Jeżeli ponadto $n^{-1}\sum_{k=1}^n\E X_k\to m$, to
dla każdego $\omega\in G$ dodajemy do powyższej
granicy granicę deterministycznych średnich i otrzymujemy
$\overline X_n(\omega)\to m$.
\end{proof}

\section{Mocne prawo wielkich liczb dla prób niezależnych jednakowo rozłożonych}

Powyższe twierdzenie zakłada zbieżność szeregu
$\sum_{n=1}^\infty\Var(X_n)/n^2$. Dla wzajemnie
niezależnych prób o tym samym rozkładzie i skończonym
drugim momencie jest to automatyczne, ponieważ ich
wariancje mają wspólną skończoną wartość $\sigma^2$
i $\sum_{n=1}^\infty n^{-2}<\infty$.
Istotne jest jednak silniejsze twierdzenie: dla
wzajemnie niezależnych zmiennych o jednakowym rozkładzie
wystarcza \emph{całkowalność}, czyli skończoność
pierwszego momentu bezwzględnego.

\begin{twierdzenie}[Mocne prawo wielkich liczb Kołmogorowa]
Niech $X_1,X_2,\ldots$ będą wzajemnie niezależnymi
rzeczywistymi zmiennymi losowymi na jednej przestrzeni
probabilistycznej $(\Omega,\mathcal F,\Prob)$,
o jednakowym rozkładzie, oraz niech $\E|X_1|<\infty$.
Wtedy
\[
\frac1n\sum_{k=1}^n X_k\longrightarrow\E X_1
\qquad\text{prawie na pewno}.
\]
\end{twierdzenie}
\begin{proof}
\emph{Krok pierwszy: obcinamy rzadkie wielkie wartości.}
Dla $k\geq1$ zdefiniujmy
$Y_k=X_k\chi_{\{|X_k|\leq k\}}$.
Każda zmienna $Y_k$ jest mierzalna jako złożenie
$X_k$ z borelowską funkcją
$g_k(x)=x\chi_{[-k,k]}(x)$.
W odróżnieniu od obcięcia przy stałym progu stosowanego
w dowodzie słabego prawa tutaj próg rośnie z indeksem.
Z jednakowości rozkładów i tożsamości warstwowej
dla nieujemnej zmiennej $|X_1|$ wynika
\[
\sum_{k=1}^\infty\Prob(|X_k|>k)
=\sum_{k=1}^\infty\Prob(|X_1|>k)
\leq\E|X_1|<\infty.
\]
Z pierwszego lematu Borela--Cantellego
(rozdział~\ref{ch:probability}) zastosowanego do
mierzalnych zdarzeń $A_k=\{|X_k|>k\}$ wynika
$\Prob(\limsup_k A_k)=0$.
Oznaczmy $G_1=\Omega\setminus\limsup_k A_k$.
Dla każdego $\omega\in G_1$ istnieje $N=N(\omega)$
takie, że dla $k\geq N$ mamy $|X_k(\omega)|\leq k$,
czyli $X_k(\omega)=Y_k(\omega)$.
Dla $n\geq N$ zachodzi zatem
\[
\frac1n\sum_{k=1}^n(X_k(\omega)-Y_k(\omega))
=\frac1n\sum_{k=1}^{N-1}(X_k(\omega)-Y_k(\omega))
\longrightarrow0.
\]
Licznik jest skończoną liczbą rzeczywistą niezależną
od $n$, a zbiór $G_1$ ma prawdopodobieństwo jeden.

\emph{Krok drugi: kontrolujemy wariancje obciętych
zmiennych.} Zmienne $Y_k$ są wzajemnie niezależne:
dla dowolnych różnych indeksów $k_1,\ldots,k_r$
i zbiorów borelowskich $B_1,\ldots,B_r$ zdarzenie
$\{Y_{k_i}\in B_i\}$ jest równe
$\{X_{k_i}\in g_{k_i}^{-1}(B_i)\}$.
Niezależność $X_k$ daje więc wymagany iloczyn
prawdopodobieństw. Nie zakładamy jednakowości
rozkładów $Y_k$, ponieważ progi obcięcia są różne.
Z $|Y_k|\leq k$ wynika $\E Y_k^2\leq k^2$,
a z definicji wariancji
$\Var(Y_k)=\E Y_k^2-(\E Y_k)^2\leq\E Y_k^2$.
Wystarczy zatem pokazać
$\sum_{k=1}^\infty\E Y_k^2/k^2<\infty$.
Z jednakowości rozkładów $X_k$ i wzoru \eqref{eq:lotus}
dla nieujemnej funkcji borelowskiej
$x\mapsto x^2\chi_{[-k,k]}(x)$ otrzymujemy
$\E Y_k^2=\E[|X_1|^2\chi_{\{|X_1|\leq k\}}]$.
Wszystkie składniki szeregu pod całką są nieujemne
i mierzalne, zatem
\[
\begin{aligned}
\sum_{k=1}^\infty\frac{\E Y_k^2}{k^2}
&=\lim_{N\to\infty}\E\left[
\sum_{k=1}^{N}\frac{|X_1|^2\chi_{\{|X_1|\leq k\}}}{k^2}
\right]\\
&=\E\left[|X_1|^2
\sum_{k=1}^\infty\frac{\chi_{\{|X_1|\leq k\}}}{k^2}\right].
\end{aligned}
\]
Ostatnia równość wynika z
\atlasref{ME-08::thm:mct}{twierdzenia o zbieżności monotonicznej};
na tym etapie dopuszczamy wartość nieskończoną.
Dla każdej liczby $x\geq0$ połóżmy
$m=\max(1,\lceil x\rceil)$, gdzie $\lceil x\rceil$
jest najmniejszą liczbą całkowitą nie mniejszą od $x$.
Wtedy zbiór dodatnich indeksów całkowitych $k$,
dla których $x\leq k$, składa się dokładnie
z indeksów $k\geq m$.
Ponieważ funkcja $t\mapsto t^{-2}$ jest malejąca
na $(0,\infty)$, mamy
\[
\sum_{k=m}^\infty\frac1{k^2}
\leq \frac1{m^2}+\int_m^\infty\frac1{t^2}\,\dd\lambda_1(t)
=\frac1{m^2}+\frac1m\leq\frac2m.
\]
Istotnie, dla każdego $k\geq m+1$ całka na
$[k-1,k]$ jest nie mniejsza niż $k^{-2}$.
Ponieważ $m\geq x$ oraz $m\geq1$, dostajemy
\[
x^2\sum_{k\geq\max(1,\lceil x\rceil)}\frac1{k^2}
\leq\frac{2x^2}{m}\leq2x
\qquad (x\geq0).
\]
Podstawiając $x=|X_1(\omega)|$ i całkując, otrzymujemy
\[
\sum_{k=1}^\infty\frac{\Var(Y_k)}{k^2}
\leq\sum_{k=1}^\infty\frac{\E Y_k^2}{k^2}
\leq2\E|X_1|<\infty.
\] Możemy więc zastosować
poprzednie twierdzenie do $Y_k$: istnieje mierzalny
zbiór $G_2$ o prawdopodobieństwie jeden, na którym
\[
\frac1n\sum_{k=1}^n(Y_k-\E Y_k)\longrightarrow0.
\]

\emph{Krok trzeci: wracamy do pierwotnych średnich.}
Z jednakowości rozkładów $X_k$ i wzoru \eqref{eq:lotus}
mamy $\E Y_k=\E[X_1\chi_{\{|X_1|\leq k\}}]$.
Na przestrzeni, na której określono $X_1$, zachodzą
punktowo $X_1\chi_{\{|X_1|\leq k\}}\to X_1$ oraz
$|X_1\chi_{\{|X_1|\leq k\}}|\leq|X_1|$.
Z założenia $|X_1|$ jest całkowalna, więc
\atlasref{ME-08::thm:dct}{twierdzenie o zbieżności zmajoryzowanej}
daje $\E Y_k\to\E X_1$.
Twierdzenie Cesàra, wykorzystane w dowodzie lematu
Kroneckera, zapewnia zatem
$n^{-1}\sum_{k=1}^n\E Y_k\to\E X_1$.
Wreszcie
\[
\frac1n\sum_{k=1}^n X_k
=\frac1n\sum_{k=1}^n(X_k-Y_k)
+\frac1n\sum_{k=1}^n(Y_k-\E Y_k)
+\frac1n\sum_{k=1}^n\E Y_k.
\]
Na mierzalnym zbiorze $G_1\cap G_2$ pierwszy i drugi
składnik dążą do zera, a trzeci (deterministyczny)
do $\E X_1$. Zbiór ten ma prawdopodobieństwo jeden,
ponieważ jego dopełnienie jest sumą dwóch zbiorów
zerowych. To kończy dowód.
\end{proof}

\begin{przyklad}[Częstości w rzutach monetą]
Niech $(B_n)$ będzie ciągiem wzajemnie niezależnych
zmiennych losowych o rozkładzie Bernoulliego
z parametrem $p\in[0,1]$, na wspólnej przestrzeni
probabilistycznej. Każda zmienna przyjmuje wartości
$0$ i $1$ prawie na pewno (wspólnie dla wszystkich $n$,
poza przeliczalną sumą zbiorów zerowych). Zmienna
$\overline B_n=n^{-1}\sum_{k=1}^n B_k$ jest częstością
sukcesów w pierwszych $n$ próbach.
Ponieważ $\E B_1=p$ i $\E|B_1|\leq1$,
MPWL daje $\overline B_n\to p$ prawie na pewno.
Twierdzenie nie stwierdza, że po dużej liczbie prób
kolejne sukcesy i porażki muszą się \emph{wyrównywać}:
dla $0<p<1$ prawdopodobieństwo samych sukcesów
w ustalonym bloku $r$ kolejnych prób wynosi $p^r>0$
z wzajemnej niezależności zmiennych.
Jeżeli do $n$ obserwacji dopiszemy $r$ kolejnych,
to
\[
|\overline B_{n+r}-\overline B_n|
=\frac{\left|\sum_{k=n+1}^{n+r}B_k-r\overline B_n\right|}
{n+r}\leq\frac{r}{n+r}
\quad\text{prawie na pewno},
\]
gdyż zarówno suma nowych obserwacji, jak i
$r\overline B_n$ należą do $[0,r]$.
Dla ustalonego $r$ wpływ takiego pasma na średnią
maleje przy $n\to\infty$.
Rysunek~\ref{fig:rp1-08-empirical-means} przedstawia trzy przykładowe przebiegi częstości sukcesów.
\end{przyklad}

\begin{figure}[htbp]
\centering
\begin{tikzpicture}[x=.041cm,y=5cm,>=Stealth]
\draw[->] (0,0)--(269,0) node[right] {$n$};
\draw[->] (0,0)--(0,1.1) node[above] {$\overline B_n$};
\foreach \x in {64,128,192,256} {
 \draw[gray!20] (\x,0)--(\x,1);
 \node[below,font=\small] at (\x,0) {$\x$};
}
\foreach \y/\lab in {0/0,.25/{1/4},.5/{1/2},.75/{3/4},1/1}
 \node[left,font=\small] at (0,\y) {$\lab$};
\draw[gray,thick,dashed] (0,.5)--(256,.5) node[right,font=\small] {$p$};
\draw[blue!75!black,thick] plot coordinates {(1,0) (2,0.5) (3,0.6666667) (4,0.5) (5,0.6) (6,0.5) (7,0.4285714) (8,0.5) (9,0.4444444) (10,0.4) (11,0.3636364) (12,0.3333333) (13,0.3076923) (14,0.2857143) (15,0.3333333) (16,0.3125) (17,0.3529412) (18,0.3888889) (19,0.4210526) (20,0.45) (21,0.4761905) (22,0.5) (23,0.4782609) (24,0.5) (25,0.52) (26,0.5) (27,0.4814815) (28,0.5) (29,0.4827586) (30,0.4666667) (31,0.4516129) (32,0.4375) (33,0.4545455) (34,0.4411765) (35,0.4285714) (36,0.4444444) (37,0.4324324) (38,0.4473684) (39,0.4358974) (40,0.45) (41,0.4634146) (42,0.452381) (43,0.4651163) (44,0.4772727) (45,0.4666667) (46,0.4565217) (47,0.4680851) (48,0.4583333) (49,0.4693878) (50,0.48) (51,0.4705882) (52,0.4615385) (53,0.4716981) (54,0.462963) (55,0.4545455) (56,0.4642857) (57,0.4736842) (58,0.4827586) (59,0.4745763) (60,0.4833333) (61,0.4754098) (62,0.4677419) (63,0.4603175) (64,0.453125) (65,0.4461538) (66,0.4545455) (67,0.4477612) (68,0.4411765) (69,0.4347826) (70,0.4285714) (71,0.4366197) (72,0.4305556) (73,0.4246575) (74,0.4324324) (75,0.44) (76,0.4473684) (77,0.4545455) (78,0.4487179) (79,0.443038) (80,0.4375) (81,0.4444444) (82,0.4512195) (83,0.4457831) (84,0.452381) (85,0.4588235) (86,0.4534884) (87,0.4482759) (88,0.4431818) (89,0.4382022) (90,0.4333333) (91,0.4285714) (92,0.4347826) (93,0.4408602) (94,0.4468085) (95,0.4526316) (96,0.4583333) (97,0.4536082) (98,0.4489796) (99,0.4444444) (100,0.44) (101,0.4356436) (102,0.4411765) (103,0.4466019) (104,0.4423077) (105,0.447619) (106,0.4528302) (107,0.4579439) (108,0.462963) (109,0.4587156) (110,0.4545455) (111,0.4504505) (112,0.4553571) (113,0.460177) (114,0.4649123) (115,0.4695652) (116,0.4655172) (117,0.4615385) (118,0.4661017) (119,0.4705882) (120,0.4666667) (121,0.4628099) (122,0.4590164) (123,0.4634146) (124,0.4677419) (125,0.464) (126,0.4603175) (127,0.4566929) (128,0.4609375) (129,0.4573643) (130,0.4538462) (131,0.4503817) (132,0.4545455) (133,0.4511278) (134,0.4552239) (135,0.4518519) (136,0.4485294) (137,0.4452555) (138,0.442029) (139,0.4460432) (140,0.45) (141,0.4468085) (142,0.443662) (143,0.4405594) (144,0.4444444) (145,0.4482759) (146,0.4452055) (147,0.4421769) (148,0.4391892) (149,0.4362416) (150,0.4333333) (151,0.4370861) (152,0.4342105) (153,0.4313725) (154,0.4350649) (155,0.4387097) (156,0.4423077) (157,0.4458599) (158,0.443038) (159,0.4402516) (160,0.44375) (161,0.447205) (162,0.4506173) (163,0.4539877) (164,0.4573171) (165,0.4606061) (166,0.4638554) (167,0.4610778) (168,0.4642857) (169,0.4615385) (170,0.4647059) (171,0.4619883) (172,0.4593023) (173,0.4624277) (174,0.4597701) (175,0.4571429) (176,0.4602273) (177,0.4576271) (178,0.4606742) (179,0.4636872) (180,0.4666667) (181,0.4696133) (182,0.467033) (183,0.4644809) (184,0.4619565) (185,0.4648649) (186,0.4623656) (187,0.459893) (188,0.4574468) (189,0.4603175) (190,0.4578947) (191,0.460733) (192,0.4635417) (193,0.4663212) (194,0.4639175) (195,0.4615385) (196,0.4591837) (197,0.4568528) (198,0.4545455) (199,0.4572864) (200,0.455) (201,0.4527363) (202,0.4554455) (203,0.453202) (204,0.4558824) (205,0.4585366) (206,0.4563107) (207,0.4589372) (208,0.4567308) (209,0.4593301) (210,0.4571429) (211,0.4597156) (212,0.4575472) (213,0.4600939) (214,0.4579439) (215,0.455814) (216,0.4537037) (217,0.4516129) (218,0.4495413) (219,0.4520548) (220,0.45) (221,0.4479638) (222,0.4459459) (223,0.4484305) (224,0.4464286) (225,0.4444444) (226,0.4469027) (227,0.4449339) (228,0.4429825) (229,0.4454148) (230,0.4434783) (231,0.4415584) (232,0.4439655) (233,0.4463519) (234,0.4444444) (235,0.4425532) (236,0.4449153) (237,0.443038) (238,0.4411765) (239,0.4393305) (240,0.4375) (241,0.4356846) (242,0.4338843) (243,0.4320988) (244,0.4303279) (245,0.4326531) (246,0.4308943) (247,0.4331984) (248,0.4354839) (249,0.4337349) (250,0.432) (251,0.4302789) (252,0.4285714) (253,0.43083) (254,0.4291339) (255,0.4313725) (256,0.4296875)};
\draw[orange!85!black,thick] plot coordinates {(1,0) (2,0.5) (3,0.6666667) (4,0.5) (5,0.6) (6,0.5) (7,0.5714286) (8,0.5) (9,0.4444444) (10,0.4) (11,0.4545455) (12,0.4166667) (13,0.3846154) (14,0.4285714) (15,0.4666667) (16,0.5) (17,0.4705882) (18,0.5) (19,0.4736842) (20,0.45) (21,0.4761905) (22,0.5) (23,0.5217391) (24,0.5416667) (25,0.56) (26,0.5384615) (27,0.5555556) (28,0.5357143) (29,0.5172414) (30,0.5333333) (31,0.5483871) (32,0.5625) (33,0.5757576) (34,0.5882353) (35,0.6) (36,0.6111111) (37,0.6216216) (38,0.6052632) (39,0.5897436) (40,0.575) (41,0.5609756) (42,0.5714286) (43,0.5813953) (44,0.5909091) (45,0.6) (46,0.6086957) (47,0.5957447) (48,0.6041667) (49,0.5918367) (50,0.6) (51,0.6078431) (52,0.6153846) (53,0.6226415) (54,0.6111111) (55,0.6) (56,0.5892857) (57,0.5789474) (58,0.5862069) (59,0.5932203) (60,0.6) (61,0.6065574) (62,0.5967742) (63,0.5873016) (64,0.59375) (65,0.5846154) (66,0.5757576) (67,0.5671642) (68,0.5588235) (69,0.5652174) (70,0.5571429) (71,0.5492958) (72,0.5555556) (73,0.5479452) (74,0.5405405) (75,0.5466667) (76,0.5394737) (77,0.5454545) (78,0.5384615) (79,0.5443038) (80,0.55) (81,0.5555556) (82,0.5609756) (83,0.5542169) (84,0.5595238) (85,0.5529412) (86,0.5581395) (87,0.5517241) (88,0.5454545) (89,0.5505618) (90,0.5444444) (91,0.5494505) (92,0.5543478) (93,0.5591398) (94,0.5531915) (95,0.5578947) (96,0.5625) (97,0.556701) (98,0.5612245) (99,0.5555556) (100,0.55) (101,0.5445545) (102,0.5392157) (103,0.5339806) (104,0.5288462) (105,0.5238095) (106,0.5188679) (107,0.5140187) (108,0.5185185) (109,0.5137615) (110,0.5181818) (111,0.5135135) (112,0.5089286) (113,0.5044248) (114,0.5087719) (115,0.5130435) (116,0.5086207) (117,0.5128205) (118,0.5084746) (119,0.5042017) (120,0.5083333) (121,0.5041322) (122,0.5) (123,0.495935) (124,0.4919355) (125,0.488) (126,0.484127) (127,0.488189) (128,0.4921875) (129,0.496124) (130,0.4923077) (131,0.4885496) (132,0.4848485) (133,0.4887218) (134,0.4925373) (135,0.4888889) (136,0.4852941) (137,0.4890511) (138,0.4855072) (139,0.4820144) (140,0.4785714) (141,0.4751773) (142,0.471831) (143,0.4685315) (144,0.4652778) (145,0.4689655) (146,0.4657534) (147,0.4693878) (148,0.4662162) (149,0.4630872) (150,0.4666667) (151,0.4635762) (152,0.4605263) (153,0.4575163) (154,0.4545455) (155,0.4580645) (156,0.4615385) (157,0.4585987) (158,0.4620253) (159,0.4591195) (160,0.45625) (161,0.4534161) (162,0.4506173) (163,0.4478528) (164,0.4512195) (165,0.4545455) (166,0.4578313) (167,0.4550898) (168,0.452381) (169,0.4556213) (170,0.4588235) (171,0.4619883) (172,0.4593023) (173,0.4624277) (174,0.4655172) (175,0.4685714) (176,0.4659091) (177,0.4689266) (178,0.4662921) (179,0.4692737) (180,0.4666667) (181,0.4640884) (182,0.467033) (183,0.4699454) (184,0.4728261) (185,0.4702703) (186,0.4731183) (187,0.4759358) (188,0.4734043) (189,0.4708995) (190,0.4736842) (191,0.4712042) (192,0.4739583) (193,0.4715026) (194,0.4742268) (195,0.4769231) (196,0.4744898) (197,0.4720812) (198,0.4747475) (199,0.4723618) (200,0.47) (201,0.4726368) (202,0.470297) (203,0.4679803) (204,0.4705882) (205,0.4731707) (206,0.4708738) (207,0.468599) (208,0.4663462) (209,0.4688995) (210,0.4714286) (211,0.4739336) (212,0.4764151) (213,0.4741784) (214,0.4766355) (215,0.4790698) (216,0.4814815) (217,0.4792627) (218,0.4770642) (219,0.4794521) (220,0.4772727) (221,0.479638) (222,0.481982) (223,0.4843049) (224,0.4866071) (225,0.4844444) (226,0.4867257) (227,0.4845815) (228,0.4824561) (229,0.4803493) (230,0.4826087) (231,0.4848485) (232,0.4827586) (233,0.4849785) (234,0.4871795) (235,0.4893617) (236,0.4872881) (237,0.4852321) (238,0.487395) (239,0.4895397) (240,0.4875) (241,0.4896266) (242,0.4876033) (243,0.4897119) (244,0.4918033) (245,0.4897959) (246,0.4878049) (247,0.48583) (248,0.483871) (249,0.4859438) (250,0.488) (251,0.4900398) (252,0.4920635) (253,0.4901186) (254,0.488189) (255,0.4901961) (256,0.4921875)};
\draw[teal!70!black,thick] plot coordinates {(1,0) (2,0) (3,0) (4,0.25) (5,0.2) (6,0.1666667) (7,0.2857143) (8,0.375) (9,0.4444444) (10,0.4) (11,0.4545455) (12,0.5) (13,0.4615385) (14,0.5) (15,0.5333333) (16,0.5625) (17,0.5294118) (18,0.5) (19,0.4736842) (20,0.45) (21,0.4761905) (22,0.4545455) (23,0.4347826) (24,0.4583333) (25,0.48) (26,0.5) (27,0.4814815) (28,0.5) (29,0.5172414) (30,0.5333333) (31,0.5483871) (32,0.5625) (33,0.5454545) (34,0.5588235) (35,0.5428571) (36,0.5555556) (37,0.5675676) (38,0.5526316) (39,0.5641026) (40,0.55) (41,0.5365854) (42,0.5238095) (43,0.5116279) (44,0.5) (45,0.5111111) (46,0.5217391) (47,0.5319149) (48,0.5416667) (49,0.5306122) (50,0.54) (51,0.5490196) (52,0.5576923) (53,0.5471698) (54,0.5555556) (55,0.5636364) (56,0.5535714) (57,0.5438596) (58,0.5344828) (59,0.5254237) (60,0.5333333) (61,0.5245902) (62,0.5322581) (63,0.5238095) (64,0.515625) (65,0.5076923) (66,0.5151515) (67,0.5223881) (68,0.5147059) (69,0.5072464) (70,0.5) (71,0.4929577) (72,0.5) (73,0.5068493) (74,0.5135135) (75,0.5066667) (76,0.5) (77,0.5064935) (78,0.5) (79,0.4936709) (80,0.4875) (81,0.4938272) (82,0.5) (83,0.4939759) (84,0.4880952) (85,0.4941176) (86,0.4883721) (87,0.4942529) (88,0.5) (89,0.505618) (90,0.5111111) (91,0.5164835) (92,0.5217391) (93,0.516129) (94,0.5106383) (95,0.5052632) (96,0.5104167) (97,0.5154639) (98,0.5204082) (99,0.5151515) (100,0.51) (101,0.5148515) (102,0.5196078) (103,0.5242718) (104,0.5192308) (105,0.5142857) (106,0.509434) (107,0.5046729) (108,0.5) (109,0.4954128) (110,0.5) (111,0.4954955) (112,0.4910714) (113,0.4867257) (114,0.4824561) (115,0.4869565) (116,0.4827586) (117,0.4871795) (118,0.4915254) (119,0.487395) (120,0.4916667) (121,0.4958678) (122,0.5) (123,0.504065) (124,0.5) (125,0.504) (126,0.5079365) (127,0.503937) (128,0.5) (129,0.496124) (130,0.4923077) (131,0.4961832) (132,0.4924242) (133,0.4887218) (134,0.4850746) (135,0.4888889) (136,0.4852941) (137,0.4817518) (138,0.4855072) (139,0.4820144) (140,0.4785714) (141,0.4822695) (142,0.4859155) (143,0.4825175) (144,0.4861111) (145,0.4827586) (146,0.4863014) (147,0.4897959) (148,0.4864865) (149,0.4832215) (150,0.48) (151,0.4834437) (152,0.4802632) (153,0.4836601) (154,0.487013) (155,0.4903226) (156,0.4935897) (157,0.4968153) (158,0.5) (159,0.5031447) (160,0.5) (161,0.5031056) (162,0.5061728) (163,0.5030675) (164,0.5060976) (165,0.5030303) (166,0.5) (167,0.502994) (168,0.5059524) (169,0.5088757) (170,0.5117647) (171,0.5087719) (172,0.5116279) (173,0.5086705) (174,0.5114943) (175,0.5085714) (176,0.5056818) (177,0.5028249) (178,0.5) (179,0.5027933) (180,0.5) (181,0.4972376) (182,0.5) (183,0.4972678) (184,0.4945652) (185,0.4918919) (186,0.4892473) (187,0.4919786) (188,0.4946809) (189,0.4920635) (190,0.4947368) (191,0.4921466) (192,0.4947917) (193,0.4974093) (194,0.5) (195,0.5025641) (196,0.5) (197,0.4974619) (198,0.4949495) (199,0.4974874) (200,0.5) (201,0.5024876) (202,0.5049505) (203,0.5024631) (204,0.504902) (205,0.5073171) (206,0.5097087) (207,0.5072464) (208,0.5048077) (209,0.5023923) (210,0.5) (211,0.4976303) (212,0.5) (213,0.4976526) (214,0.4953271) (215,0.4930233) (216,0.4953704) (217,0.4930876) (218,0.4954128) (219,0.4931507) (220,0.4954545) (221,0.4977376) (222,0.4954955) (223,0.4932735) (224,0.4910714) (225,0.4888889) (226,0.4867257) (227,0.4845815) (228,0.4868421) (229,0.4847162) (230,0.4826087) (231,0.4848485) (232,0.4827586) (233,0.4806867) (234,0.4786325) (235,0.4808511) (236,0.4788136) (237,0.4810127) (238,0.4831933) (239,0.4853556) (240,0.4833333) (241,0.4813278) (242,0.4793388) (243,0.4773663) (244,0.4754098) (245,0.477551) (246,0.4756098) (247,0.4736842) (248,0.4758065) (249,0.4779116) (250,0.476) (251,0.4741036) (252,0.4761905) (253,0.4782609) (254,0.480315) (255,0.4784314) (256,0.4765625)};
\end{tikzpicture}
\caption{Trzy przykładowe przebiegi średniej empirycznej
dla binarnych ciągów wygenerowanych pseudolosowo
z parametrem $p=1/2$. Każdy punkt jest średnią
pierwszych $n$ obserwacji z danego ciągu.
Wykres ilustruje mocne prawo wielkich liczb,
ale skończona symulacja nie stanowi jego dowodu.}
\label{fig:rp1-08-empirical-means}
\end{figure}


\begin{przyklad}[Dlaczego dla rozkładu Cauchy'ego nie ma analogicznego MPWL]
Niech $C_1,C_2,\ldots$ będą wzajemnie niezależnymi
zmiennymi losowymi na jednej przestrzeni probabilistycznej
i niech każda ma standardowy rozkład Cauchy'ego.
Jego pierwszy moment bezwzględny jest nieskończony
(rozdział~\ref{ch:distributions}), więc hipoteza
całkowalności w MPWL Kołmogorowa nie jest spełniona.
Z wzoru na funkcję charakterystyczną standardowego
rozkładu Cauchy'ego $\varphi_C(t)=e^{-|t|}$ oraz
faktoryzacji funkcji charakterystycznych sumy zmiennych
wzajemnie niezależnych (rozdział~\ref{ch:expectation})
wynika dla każdego całkowitego $n\geq1$ i $t\in\R$,
że
\[
\varphi_{(C_1+\cdots+C_n)/n}(t)
=\bigl[\varphi_C(t/n)\bigr]^n=e^{-|t|}.
\]
Z twierdzenia o jednoznaczności rozkładu wyznaczonego
przez funkcję charakterystyczną
(rozdział~\ref{ch:expectation}) wynika, że każda średnia
$n^{-1}\sum_{k=1}^n C_k$ ma ten sam rozkład co $C_1$.
Gdyby dla pewnego $c\in\R$ średnie zbiegały do $c$
według prawdopodobieństwa, to z twierdzenia
o podstawowych implikacjach zbieżności
(rozdział~\ref{ch:convergence}) zbiegałyby do $c$
również według rozkładu. Jednocześnie ich rozkłady
tworzą stały ciąg miar o granicy $\Law(C_1)$.
Z jednoznaczności granicy według rozkładu wynikałoby
$\Law(C_1)=\delta_c$. Jest to niemożliwe, ponieważ
rozkład Cauchy'ego ma gęstość względem miary Lebesgue'a,
a zatem $\Prob(C_1=c)=0$, podczas gdy $\delta_c(\{c\})=1$.
Nie zachodzi więc ani SPWL, ani MPWL z granicą stałą.
Brak całkowalności nie jest tylko techniczną luką
dowodu prawa wielkich liczb.
\end{przyklad}

\section{Alternatywny dowód MPWL: kontrola podciągu diadycznego}

Dla wzajemnie niezależnych zmiennych o wspólnej średniej
i jednostajnie ograniczonych wariancjach można udowodnić
mocne prawo także bez lematu Kroneckera. Poniższy dowód
pokazuje, jak kontrola \emph{maksimum} sum częściowych
pozwala przejść od oszacowań na indeksach diadycznych
do zbieżności całego ciągu prawie na pewno.

\begin{twierdzenie}[MPWL dla niezależnych zmiennych o jednostajnie ograniczonych wariancjach]
Niech $X_1,X_2,\ldots$ będą wzajemnie niezależnymi
rzeczywistymi zmiennymi losowymi na jednej przestrzeni
probabilistycznej $(\Omega,\mathcal F,\Prob)$,
o skończonych drugich momentach, wspólnej średniej
$m\in\R$ oraz $\sup_n\Var(X_n)\leq C<\infty$.
Wówczas $n^{-1}\sum_{k=1}^nX_k\to m$ prawie na pewno.
\end{twierdzenie}
\begin{proof}
Połóżmy $Z_k=X_k-m$ i $S_n=\sum_{k=1}^nZ_k$.
Zmienne $Z_k$ są wzajemnie niezależne, ponieważ
powstają z $X_k$ przez osobne przekształcenia borelowskie.
Mają średnie zero oraz
$\Var(Z_k)=\Var(X_k)\leq C$, a więc spełniają hipotezy
nierówności maksymalnej Kołmogorowa.
Dla $\varepsilon>0$ oraz całkowitego $j\geq1$ określmy
mierzalne zdarzenie
$A_j(\varepsilon)=\{\max_{1\leq k\leq2^j}|S_k|>
\varepsilon2^j\}$.
Z nierówności maksymalnej otrzymujemy
\[
\Prob\left(\max_{1\leq k\leq2^j}|S_k|>
\varepsilon2^j\right)
\leq\frac{\sum_{k\leq2^j}\Var(Z_k)}
{\varepsilon^2 2^{2j}}
\leq\frac{C}{\varepsilon^2 2^j}.
\]
Ponieważ $\sum_{j=1}^\infty2^{-j}=1$, mamy
$\sum_{j=1}^\infty\Prob(A_j(\varepsilon))<\infty$.
Z pierwszego lematu Borela--Cantellego
(rozdział~\ref{ch:probability}) wynika, że mierzalny
zbiór $G_\varepsilon=\Omega\setminus\limsup_j A_j(\varepsilon)$
ma prawdopodobieństwo jeden.
Dla każdego $\omega\in G_\varepsilon$ istnieje
$J=J(\omega,\varepsilon)$ takie, że dla $j\geq J$
oraz $1\leq k\leq 2^j$ mamy $|S_k(\omega)|\leq\varepsilon 2^j$.
Każdemu całkowitemu $n\geq2$ odpowiada dokładnie
jeden całkowity $j\geq1$ spełniający $2^{j-1}<n\leq2^j$,
a ten indeks dąży do nieskończoności wraz z $n$.
Dla dostatecznie dużego $n$ otrzymujemy
\[
\frac{|S_n(\omega)|}{n}
\leq\frac{\varepsilon2^j}{n}<2\varepsilon.
\]
Mierzalny zbiór $G=\bigcap_{r=1}^\infty G_{1/r}$
ma prawdopodobieństwo jeden, gdyż jego dopełnienie
jest przeliczalną sumą zbiorów zerowych.
Dla każdego $\omega\in G$ oraz każdego $r\geq1$
mamy $\limsup_{n\to\infty}|S_n(\omega)|/n\leq2/r$.
Przechodząc z $r\to\infty$, otrzymujemy $S_n(\omega)/n\to0$.
Ponieważ $\overline X_n=m+S_n/n$, zachodzi teza.
\end{proof}

\begin{uwaga}
Przy zachowaniu wzajemnej niezależności założenie
jednakowego rozkładu nie jest tu potrzebne: wystarczy
wspólna średnia i jednostajne ograniczenie wariancji.
Twierdzenie jest szczególnym przypadkiem wcześniejszego
MPWL przy sumowalnych wariancjach ważonych, ponieważ
$\sum_{n=1}^\infty\Var(X_n)/n^2
\leq C\sum_{n=1}^\infty n^{-2}<\infty$.
Natomiast w przypadku jednakowych rozkładów MPWL
Kołmogorowa wymaga tylko całkowalności, a nie
skończoności drugiego momentu. W ogólności tych dwóch
zestawów założeń (jednakowych rozkładów i całkowalności
oraz wspólnej średniej i ograniczonych wariancji)
nie należy utożsamiać.
\end{uwaga}

\begin{przyklad}[MPWL pomimo nieskończonej wariancji]
Niech $X_1,X_2,\ldots$ będą wzajemnie niezależnymi
zmiennymi losowymi na wspólnej przestrzeni
probabilistycznej, o jednakowym rozkładzie Pareto
z gęstością
\[
f(x)=
\begin{cases}
\alpha x^{-\alpha-1},&x\geq1,\\
0,&x<1,
\end{cases}
\qquad 1<\alpha<2.
\]
Jest to nieujemna funkcja borelowska, a
$\int_1^\infty\alpha x^{-\alpha-1}\,\dd\lambda_1(x)=1$,
ponieważ całka na $[1,R]$ wynosi $1-R^{-\alpha}$.
Zatem $f$ jest gęstością probabilistyczną.
Ponieważ rozkład jest skupiony na $[1,\infty)$,
mamy $X_1\geq1$ prawie na pewno. Ze wzoru
\eqref{eq:lotus} dla funkcji nieujemnych oraz
\atlasref{ME-08::thm:mct}{twierdzenia o zbieżności monotonicznej}
dla całek obciętych do $[1,R]$ (najpierw dla
całkowitych $R\to\infty$) otrzymujemy
\[
\begin{aligned}
\E|X_1|=\E X_1
&=\lim_{R\to\infty}\alpha\int_1^R x^{-\alpha}\,
\dd\lambda_1(x)=\frac{\alpha}{\alpha-1},\\
\E X_1^2
&=\lim_{R\to\infty}\alpha\int_1^R x^{1-\alpha}\,
\dd\lambda_1(x)\\
&=\lim_{R\to\infty}\frac{\alpha}{2-\alpha}
(R^{2-\alpha}-1)=+\infty.
\end{aligned}
\]
W pierwszej równości granicznej korzystamy z
$\int_1^R x^{-\alpha}\,\dd\lambda_1(x)
=(1-R^{1-\alpha})/(\alpha-1)$.
Jeśli $m=\E X_1$, to dla $x\geq 2m$ mamy
$(x-m)^2\geq x^2/4$; całka drugiego momentu na
$[2m,\infty)$ jest nieskończona. Zatem również
$\E(X_1-m)^2=+\infty$.
W tym sensie mówimy o nieskończonej wariancji:
całkę definiującą wariancję dopuszczamy tutaj
w rozszerzonym przedziale $[0,\infty]$, mimo że
w algebrze momentów rozważaliśmy wariancje skończone
(rozdział~\ref{ch:expectation}).
Wariant MPWL wykorzystujący sumowalne wariancje
ważone nie ma tu zastosowania, ale MPWL Kołmogorowa
zapewnia $\overline X_n\to\alpha/(\alpha-1)$ p.n.
\end{przyklad}

\begin{przyklad}[Dlaczego nie można pominąć warunku niezależności]
Niech $Z$ ma rozkład Bernoulliego z parametrem $1/2$
i połóżmy $X_n=Z$ dla każdego $n$. Wszystkie $X_n$
mają ten sam rozkład i skończone momenty dowolnego
rzędu, lecz $\overline X_n=Z$ dla wszystkich $n$.
Dla $j\ne k$ mamy
$\Prob(X_j=1,X_k=1)=1/2\ne1/4
=\Prob(X_j=1)\Prob(X_k=1)$, więc zmienne nie są
nawet parami niezależne.
Dla każdego $n$ zachodzi
$\Prob(|\overline X_n-1/2|>1/4)=1$,
a zatem nie ma zbieżności według prawdopodobieństwa
do $\E Z=1/2$. Granicą prawie na pewno jest natomiast
$Z$, które różni się od $1/2$ z prawdopodobieństwem jeden.
Kolejne obserwacje nie wnoszą nowej informacji,
jeśli są identyczną kopią tej samej zmiennej.
\end{przyklad}

\begin{przyklad}[Różnica między zbieżnością średnich a zbieżnością pojedynczych wyników]
Dla wzajemnie niezależnych zmiennych
$B_n\sim\operatorname{Bernoulli}(1/2)$ na wspólnej
przestrzeni probabilistycznej MPWL Kołmogorowa daje
$\overline B_n\to1/2$ p.n.
Rodzina zdarzeń $D_n=\{B_n=1\}$ jest wzajemnie niezależna,
podobnie jak rodzina $E_n=\{B_n=0\}$, z definicji
niezależności zmiennych losowych. Ponieważ
$\Prob(D_n)=\Prob(E_n)=1/2$, oba szeregi
$\sum_n\Prob(D_n)$ i $\sum_n\Prob(E_n)$ są rozbieżne.
Drugi lemat Borela--Cantellego (rozdział~\ref{ch:probability})
daje $\Prob(\limsup_n D_n)=\Prob(\limsup_n E_n)=1$.
Na przecięciu tych dwóch zbiorów pełnej miary
ciąg $B_n(\omega)$ zawiera podciąg stale równy jeden
i podciąg stale równy zero. Z jednoznaczności granicy
liczbowej wynika, że ciąg $B_n$ prawie na pewno
\emph{nie jest} zbieżny.
To stabilizacja średniej, nie stabilizacja
poszczególnych rezultatów.
\end{przyklad}

